(1) 作用量为 \(\int dt\,dA\{\rho h\,w_t^2/2-[D_xw_{xx}^2+2D_{xy}w_{xx}w_{yy}+D_yw_{yy}^2]/2\}\)。对 \(w\) 变分并两次分部积分，在无限板无边界项条件下得 \(\boxed{\rho h\,w_{tt}+D_xw_{xxxx}+2D_{xy}w_{xxyy}+D_yw_{yyyy}=0}\)。交叉项两次变分贡献系数2，不能错写成 \(D_{xy}\)；方程描述弯曲主导的四阶空间波，不是张紧膜的二阶波动方程。

(2) 代入平面波得到 \(\boxed{\omega^2=[D_xk_x^4+2D_{xy}k_x^2k_y^2+D_yk_y^4]/(\rho h)}\)。令 \(S=D_xk_x^4+2D_{xy}k_x^2k_y^2+D_yk_y^4\)，则 \(v_{gx}=\partial_{k_x}\omega=2(D_xk_x^3+D_{xy}k_xk_y^2)/(\rho h\,\omega)\)，\(v_{gy}=2(D_yk_y^3+D_{xy}k_x^2k_y)/(\rho h\,\omega)\)。例如 \(k_x=k_y>0\) 且 \(D_x\ne D_y\) 时两分量不相等，故群速度不沿 \(45^\circ\) 波矢方向；各向异性会导致能量射线偏折。

(3) 简支边界允许 \(w=0\) 且相应弯矩为零，正弦双指标函数满足这些条件。取 \(k_x=m\pi/a,k_y=n\pi/b\) 得 \(\boxed{\omega_{mn}^2=[D_x(m\pi/a)^4+2D_{xy}(m\pi/a)^2(n\pi/b)^2+D_y(n\pi/b)^4]/(\rho h)}\)。利用 \(\int_0^a\sin^2(m\pi x/a)dx=a/2\) 及 \(y\) 方向同式，峰值势能为 \(E_{mn}=abA^2S_{mn}/8=\rho habA^2\omega_{mn}^2/8\)。峰值动能相等，周期中二者互换，因无耗散总能量不变。

(4) 二次型 \(D_xX^2+2D_{xy}XY+D_yY^2\) 在题设 \(|D_{xy}|<\sqrt{D_xD_y}\) 下严格正定，保证所有非零波矢频率实。在形式各向同性色散极限 \(D_x=D_y=D_{xy}=D\) 时 \(S=D(k_x^2+k_y^2)^2\)；该极限处于 \(|D_{xy}|<\sqrt{D_xD_y}\) 的边界，若要描述完整各向同性板的所有曲率，还须保留本题忽略的扭转曲率能项，故 \(\omega=\sqrt{D/(\rho h)}\,k^2\)，\(v_g=2\sqrt{D/(\rho h)}\,k\hat k\)，相速为其一半。固定传播方向令 \(k\to0\)，二者都正比 \(k\) 而趋零；当波长与板厚同阶，Kirchhoff–Love忽略的横向剪切与转动惯量不可忽略。
